\documentclass{plantilla}

\usepackage{siunitx}
\usepackage{array}
\usepackage{ifthen}
% Ningun flotante puede quedar antes del lugar donde se lo declara
\usepackage{flafter}
\usepgfplotslibrary{groupplots}
\usetikzlibrary{positioning, arrows.meta, fit, calc}

% Mas espacio para flotantes en cada columna
\renewcommand{\topfraction}{0.9}
\renewcommand{\bottomfraction}{0.7}
\renewcommand{\textfraction}{0.1}
\renewcommand{\floatpagefraction}{0.85}
\renewcommand{\dbltopfraction}{0.9}
\setcounter{topnumber}{3}
\setcounter{bottomnumber}{2}
\setcounter{totalnumber}{5}

\sisetup{
    output-decimal-marker = {,},
    group-digits = none,
    exponent-product = \cdot,
    uncertainty-mode = separate,
    per-mode = symbol
}

% Numeros de las mediciones y de la simulacion (generados por generar_datos.py y generar_ej2.py)
% Generado por generar_datos.py: no editar a mano
\newcommand{\resNPuntos}{17}
\newcommand{\resNBuenos}{15}
\newcommand{\resVinPP}{1,13}
\newcommand{\resOffsetBarrido}{0,72}
\newcommand{\resFTres}{9,98}
\newcommand{\resDbCincuenta}{-28,10}
\newcommand{\resPend}{-41,4}
\newcommand{\resFcFit}{10,0000}
\newcommand{\resSFcFit}{0,0004}
\newcommand{\resQFit}{0,7071}
\newcommand{\resSQFit}{0,0001}
\newcommand{\resKFit}{1,0001}
\newcommand{\resSKFit}{0,0001}
\newcommand{\resMaxEdb}{0,005}
\newcommand{\resMaxEfa}{0,03}
\newcommand{\resDtDiez}{25,0}
\newcommand{\resMp}{4,34}
\newcommand{\resTp}{70}
\newcommand{\resTs}{95}
\newcommand{\resZeta}{0,707}
\newcommand{\resWn}{63,4}
\newcommand{\resFn}{10,10}
\newcommand{\resQEsc}{0,707}
\newcommand{\resMpSim}{4,33}
\newcommand{\resTpSim}{70}
\newcommand{\resTsSim}{95}
\newcommand{\resEscBajo}{0,65}
\newcommand{\resEscAlto}{3,04}
\newcommand{\resFcId}{10,10}
\newcommand{\resQId}{0,707}
\newcommand{\resBcero}{9.6253e-04}
\newcommand{\resBuno}{1.9251e-03}
\newcommand{\resBdos}{9.6253e-04}
\newcommand{\resAuno}{-1,91039}
\newcommand{\resAdos}{0,91424}
\newcommand{\resFoSK}{9,68}
\newcommand{\resQSK}{0,709}
\newcommand{\resFtresSK}{9,69}
\newcommand{\resDbCincuentaSK}{-28,5}
\newcommand{\resDbCincuentaZ}{-28,10}
\newcommand{\resTPdosFc}{10,83}
\newcommand{\resTPdosDb}{-27,8}
\newcommand{\resGanDC}{1,0000}
\newcommand{\resOrdDC}{0,0}
\newcommand{\resErrMax}{0,2}

% Generado por generar_ej2.py: no editar a mano
\newcommand{\ejFs}{200}
\newcommand{\ejFpuno}{27163}
\newcommand{\ejFpdos}{28863}
\newcommand{\ejEps}{0,15262}
\newcommand{\ejQ}{2,5137}
\newcommand{\ejGdb}{19,08}
\newcommand{\ejSpec}{18,98}
\newcommand{\ejKbil}{3.73869e+05}
\newcommand{\ejBfo}{1,19603}
\newcommand{\ejAfo}{-1,10543}
\newcommand{\ejAdfo}{0,73421}
\newcommand{\ejBb}{1,40238}
\newcommand{\ejAb}{-1,07606}
\newcommand{\ejAdb}{0,68836}
\newcommand{\ejFoB}{28004}
\newcommand{\ejQB}{2,0874}
\newcommand{\ejDbFoUno}{-0,129}
\newcommand{\ejDbFoDos}{-0,131}
\newcommand{\ejDbBUno}{-0,090}
\newcommand{\ejDbBDos}{-0,090}
\newcommand{\ejDbSUno}{-0,100}
\newcommand{\ejDbSDos}{-0,100}
\newcommand{\ejDbCentroB}{19,08}
\newcommand{\ejFactor}{14,2}
\newcommand{\ejTresSuno}{22,99}
\newcommand{\ejTresSdos}{34,10}
\newcommand{\ejTresFouno}{23,47}
\newcommand{\ejTresFodos}{33,13}
\newcommand{\ejTresBuno}{22,64}
\newcommand{\ejTresBdos}{34,23}
\newcommand{\ejDbNoventaZ}{-9,8}
\newcommand{\ejDbNoventaS}{1,7}
\newcommand{\ejRFo}{0,8569}
\newcommand{\ejRB}{0,8297}
\newcommand{\ejRFormula}{0,8569}
\newcommand{\ejAngB}{49,57}
\newcommand{\ejNuno}{25}
\newcommand{\ejTUno}{125}
\newcommand{\ejHcero}{1,402}
\newcommand{\ejMinBits}{12}
\newcommand{\ejCiclos}{360}
\newcommand{\ejFsMaxPlaca}{77}
\newcommand{\ejTonoUno}{14}
\newcommand{\ejTonoUnoMed}{8,09}
\newcommand{\ejTonoUnoTeo}{8,09}
\newcommand{\ejTonoDos}{28}
\newcommand{\ejTonoDosMed}{19,08}
\newcommand{\ejTonoDosTeo}{19,08}
\newcommand{\ejTonoTres}{56}
\newcommand{\ejTonoTresMed}{5,72}
\newcommand{\ejTonoTresTeo}{5,72}

% \val{numero}{unidad} y \nm{numero}: como \SI y \num, pero muestran ?? si el dato falta
\newcommand{\val}[2]{\ifthenelse{\equal{#1}{??}}{\textbf{??}\,\si{#2}}{\SI{#1}{#2}}}
\newcommand{\nm}[1]{\ifthenelse{\equal{#1}{??}}{\textbf{??}}{\num{#1}}}

\materia{Teoría de Circuitos II}
\comision{4R2}
\titulo{Trabajo Práctico N\textsuperscript{o} 3: Filtros Digitales}
\autores{Grupo 2: Ernst Pedro (400624), Rao Valentino (402308), Gil Ignacio (401891)}

\begin{document}

\maketitle

\begin{abstract}
En el Ejercicio~1 se identifica el filtro IIR grabado en una Blue Pill (STM32F103C8), que debe reemplazar al Sallen-Key del módulo de \SIrange{4}{20}{\milli\ampere} del TP2: pasa bajos Butterworth de segundo orden, $f_c = \SI{10}{\hertz}$ y ganancia unitaria. A partir de la caracterización estática, la respuesta al escalón y la respuesta en frecuencia se obtienen $H(s)$ y $H(z)$, se las compara con el Sallen-Key y se muestra el efecto del aliasing. En el Ejercicio~2 se diseña, por transformación bilineal con predistorsión, el equivalente digital del pasa banda del TP2 ($f_0 = \SI{28}{\kilo\hertz}$, $BW = \SI{1,7}{\kilo\hertz}$, $K = 9$) a $f_s = \SI{200}{\kilo\hertz}$. Para cumplir la especificación hay que predistorsionar los bordes de banda, y los coeficientes necesitan al menos \ejMinBits{} bits.
\end{abstract}

\begin{IEEEkeywords}
Filtro digital IIR, transformación bilineal, predistorsión, aliasing, cuantización, STM32.
\end{IEEEkeywords}

\section{Introducción}

En el Ejercicio~1, una Blue Pill~\cite{ds} muestrea a $f_s = \SI{1}{\kilo\hertz}$ la señal de un generador aplicada en PA0, la filtra con un IIR ya grabado y envía a la PC bloques de 2048 muestras de la entrada y la salida. Ese filtro debería reemplazar a la Etapa~2 (Sallen-Key) del módulo de \SIrange{4}{20}{\milli\ampere} del TP2~\cite{tp2}, y la tarea es verificar con mediciones si cumple la especificación y obtener su $H(s)$ y su $H(z)$. La corriente de lazo se representa con su equivalente a la entrada del ADC, \SI{0,6}{\volt} para \SI{4}{\milli\ampere} y \SI{3,0}{\volt} para \SI{20}{\milli\ampere} (\SI{0,15}{\volt\per\milli\ampere}). En el Ejercicio~2 se diseña, solo por simulación, el filtro digital equivalente al filtro del TP2 que no se implementó físicamente.

\section{Ejercicio 1: filtro de la señal de 4--20 mA}

\subsection{Montaje y método}
El generador se conectó directamente a PA0, sin la resistencia de protección de \SI{1}{\kilo\ohm}, que no altera la medición mientras la tensión se mantenga entre 0 y \SI{3,3}{\volt}. Como no se dispuso de osciloscopio ni multímetro, los niveles se verificaron con la propia lectura del ADC, acercándose a cada valor desde abajo para no superar ese rango. Con el \texttt{graficador.py} provisto se perdía cerca de un tercio de cada bloque y casi siempre la línea final \texttt{\#FIN}, porque leía el puerto de a un byte y no lo vaciaba al ritmo de las ráfagas de la placa. Se lo corrigió para que lea en bloques, desde un proceso separado de la interfaz gráfica, y acepte solo bloques completos. Además se le agregó la medición en vivo de ganancia y fase por ajuste senoidal de mínimos cuadrados sobre las 2048 muestras~\cite{ieee1057}, y una tecla que registra cada punto de un barrido en un archivo CSV. Los niveles de continua se toman como el valor medio del bloque.

\subsection{Caracterización estática}
La Tabla~\ref{tab:estatica} muestra la entrada y la salida medias para los cinco niveles aplicados, cercanos a 4, 8, 12, 16 y \SI{20}{\milli\ampere} y fijados con un generador digital en continua. $I_i$ e $I_o$ son las corrientes equivalentes, y el error es $I_o - I_i = (V_o - V_i)/(\SI{0,15}{\volt\per\milli\ampere})$. La recta de ajuste de $V_o$ en función de $V_i$ tiene pendiente (ganancia de continua) \nm{\resGanDC} y ordenada \val{\resOrdDC}{\milli\volt}, y el error máximo es de \val{\resErrMax}{\micro\ampere}. Es decir, la ganancia de continua es unitaria, como pide la especificación, y el error queda muy por debajo del escalón del ADC (\SI{0,8}{\milli\volt}, unos \SI{5}{\micro\ampere}).

\begin{table}[!htb]
\centering
\caption{Caracterización estática.}
\label{tab:estatica}
\begin{tabular}{ccccc}
\toprule
$I_i$ (mA) & $V_i$ (V) & $V_o$ (V) & $I_o$ (mA) & Error (\si{\micro\ampere}) \\
\midrule
\num{4,100} & \num{0,6151} & \num{0,6151} & \num{4,100} & \num{0,0} \\
\num{8,013} & \num{1,2020} & \num{1,2021} & \num{8,014} & \num{0,2} \\
\num{12,005} & \num{1,8008} & \num{1,8007} & \num{12,005} & \num{-0,2} \\
\num{15,987} & \num{2,3981} & \num{2,3981} & \num{15,987} & \num{0,0} \\
\num{19,662} & \num{2,9492} & \num{2,9492} & \num{19,662} & \num{0,0} \\
\bottomrule

\end{tabular}
\end{table}

\subsection{Respuesta al escalón}
Se aplicó una onda cuadrada de \SI{0,5}{\hertz} entre \val{\resEscBajo}{\volt} y \val{\resEscAlto}{\volt} (Fig.~\ref{fig:escalon}). Se midió un sobrepico $M_p = \val{\resMp}{\percent}$, un tiempo de pico $t_p = \val{\resTp}{\milli\second}$ y un tiempo de establecimiento al \SI{2}{\percent} de $t_s = \val{\resTs}{\milli\second}$, con la resolución de \SI{1}{\milli\second} que da el muestreo. Para un sistema de segundo orden,
\begin{equation}
\zeta = \frac{-\ln M_p}{\sqrt{\pi^2 + \ln^2 M_p}}, \qquad \omega_n = \frac{\pi}{t_p\sqrt{1 - \zeta^2}},
\label{eq:zetawn}
\end{equation}
lo que da $\zeta = \nm{\resZeta}$ y $\omega_n = \val{\resWn}{\radian\per\second}$ ($f_n = \val{\resFn}{\hertz}$). La función de transferencia propuesta es entonces
\begin{equation}
H(s) = \frac{\omega_n^2}{s^2 + 2\zeta\omega_n s + \omega_n^2}.
\label{eq:hs}
\end{equation}
El Butterworth nominal ($\zeta = 0{,}707$, \SI{10}{\hertz}) discretizado da $M_p = \SI{\resMpSim}{\percent}$, $t_p = \SI{\resTpSim}{\milli\second}$ y $t_s = \SI{\resTsSim}{\milli\second}$, prácticamente lo mismo que se midió. El $\zeta$ medido es el de Butterworth. La diferencia de $f_n$ con \SI{10}{\hertz} se debe a la resolución de $t_p$: el valor teórico es \SI{70,7}{\milli\second}, pero solo puede medirse en múltiplos de \SI{1}{\milli\second}.

\begin{figure}[!htb]
\centering
\begin{tikzpicture}
\begin{axis}[width=\columnwidth, height=4.2cm, xmin=-50, xmax=400, ymin=0.3, ymax=3.4,
    xlabel={$t$ (ms)}, ylabel={Tensión (V)},
    grid=both, grid style={line width=0.1pt, draw=gray!30}, major grid style={line width=0.2pt, draw=gray!50},
    tick label style={font=\footnotesize}, label style={font=\footnotesize},
    legend style={font=\scriptsize, cells={anchor=west}, at={(0.98,0.05)}, anchor=south east},
    /pgf/number format/use comma, unbounded coords=jump]
\addplot[gray, thin] table[x=t, y=x] {datos/escalon.dat};
\addlegendentry{Entrada}
\addplot[only marks, mark=*, mark size=0.6pt, blue] table[x=t, y=y] {datos/escalon.dat};
\addlegendentry{Salida medida}
\addplot[red, thin] table[x=t, y=sim] {datos/escalon.dat};
\addlegendentry{$H(z)$ nominal}
\end{axis}
\end{tikzpicture}
\caption{Respuesta al escalón de \protect\val{\resEscBajo}{\volt} a \protect\val{\resEscAlto}{\volt}, medida y simulada con el Butterworth nominal discretizado.}
\label{fig:escalon}
\end{figure}

\subsection{Respuesta en frecuencia}
Se barrió de 1 a \SI{100}{\hertz} con una senoidal de \val{\resVinPP}{\volt} pico a pico montada sobre \val{\resOffsetBarrido}{\volt} de continua (Tabla~\ref{tab:barrido} y Fig.~\ref{fig:bode}). Como el filtro es lineal, el nivel de continua no modifica la ganancia ni la fase. El corrimiento temporal entre entrada y salida es $\Delta t = -\varphi/(360^\circ f)$: por ejemplo, a \SI{10}{\hertz} la salida se atrasa \val{\resDtDiez}{\milli\second}. La ganancia cae \SI{3}{\dB} en \val{\resFTres}{\hertz}, la atenuación a \SI{50}{\hertz} es de \val{\resDbCincuenta}{\dB}, y entre 50 y \SI{100}{\hertz} la pendiente es de \val{\resPend}{\dB} por década, la de un filtro de orden~2. Ajustando $f_c$, $Q$ y $K$ del modelo discretizado a todos los puntos, salvo dos tomados mientras se cambiaba la frecuencia ($^{*}$), resulta $f_c = \val{\resFcFit \pm \resSFcFit}{\hertz}$, $Q = \nm{\resQFit \pm \resSQFit}$ y $K = \nm{\resKFit \pm \resSKFit}$. Respecto del Butterworth nominal discretizado, el error máximo es de \val{\resMaxEdb}{\dB} y \val{\resMaxEfa}{\degree}: el filtro cumple la especificación.

\begin{table}[!htb]
\centering
\caption{Barrido en frecuencia y diferencia con el Butterworth nominal discretizado.}
\label{tab:barrido}
\footnotesize
\setlength{\tabcolsep}{3.5pt}
\begin{tabular}{rrrrrr}
\toprule
$f$ (Hz) & $|H|$ (dB) & $\varphi$ (°) & $\Delta t$ (ms) & Dif. (dB) & Dif. (°) \\
\midrule
\num{1,00} & \num{0,00} & \num{-8,1} & \num{22,5} & \num{0,00} & \num{0,0} \\
\num{2,00} & \num{-0,01} & \num{-16,4} & \num{22,8} & \num{0,00} & \num{0,0} \\
\num{3,00} & \num{-0,04} & \num{-25,0} & \num{23,1} & \num{0,00} & \num{0,0} \\
\num{5,00} & \num{-0,26} & \num{-43,3} & \num{24,1} & \num{0,00} & \num{0,0} \\
\num{5,00} & \num{-0,26} & \num{-43,3} & \num{24,1} & \num{0,00} & \num{0,0} \\
\num{7,00} & \num{-0,93} & \num{-62,7} & \num{24,9} & \num{0,00} & \num{0,0} \\
\num{8,00} & \num{-1,49} & \num{-72,3} & \num{25,1} & \num{0,00} & \num{0,0} \\
\num{8,99}$^{*}$ & \num{-2,22} & \num{-80,9} & \num{25,0} & \num{-0,04} & \num{0,5} \\
\num{10,00} & \num{-3,01} & \num{-90,0} & \num{25,0} & \num{0,00} & \num{0,0} \\
\num{11,00} & \num{-3,92} & \num{-97,7} & \num{24,7} & \num{0,00} & \num{0,0} \\
\num{11,94}$^{*}$ & \num{-4,91} & \num{-103,1} & \num{24,0} & \num{-0,09} & \num{1,1} \\
\num{15,00} & \num{-7,83} & \num{-120,5} & \num{22,3} & \num{0,00} & \num{0,0} \\
\num{20,00} & \num{-12,32} & \num{-136,7} & \num{19,0} & \num{0,00} & \num{0,0} \\
\num{30,00} & \num{-19,18} & \num{-152,1} & \num{14,1} & \num{0,00} & \num{0,0} \\
\num{49,99} & \num{-28,10} & \num{-163,7} & \num{9,1} & \num{0,00} & \num{0,0} \\
\num{69,99} & \num{-34,08} & \num{-168,6} & \num{6,7} & \num{0,00} & \num{0,0} \\
\num{99,99} & \num{-40,57} & \num{-172,1} & \num{4,8} & \num{0,00} & \num{0,0} \\
\bottomrule

\end{tabular}
\end{table}

\subsection{Función de transferencia discreta}
Con la $H(s)$ de~\eqref{eq:hs}, la transformación bilineal $s = k(1 - z^{-1})/(1 + z^{-1})$ con predistorsión en $f_c$, $k = 2\pi f_c/\tan(\pi f_c/f_s)$~\cite{oppenheim}, da
\begin{equation}
H(z) = \frac{b_0 (1 + 2z^{-1} + z^{-2})}{1 + a_1 z^{-1} + a_2 z^{-2}},
\label{eq:hz}
\end{equation}
con $b_0 = \num{\resBcero}$, $a_1 = \num{\resAuno}$ y $a_2 = \num{\resAdos}$ (para $f_c = \SI{\resFcId}{\hertz}$, $Q = \num{\resQId}$). La Fig.~\ref{fig:bode} superpone $|H(e^{j\omega})|$ con los puntos medidos.

\begin{figure}[!htb]
\centering
\begin{tikzpicture}
\begin{groupplot}[
    group style={group size=1 by 2, vertical sep=3mm, xticklabels at=edge bottom},
    width=\columnwidth, height=4.0cm, xmode=log, xmin=0.8, xmax=130,
    grid=both, grid style={line width=0.1pt, draw=gray!30}, major grid style={line width=0.2pt, draw=gray!50},
    tick label style={font=\footnotesize}, label style={font=\footnotesize},
    legend style={font=\scriptsize, cells={anchor=west}}, /pgf/number format/use comma]
\nextgroupplot[ylabel={$|H|$ (dB)}, ymin=-45, ymax=3, ytick={-40,-30,-20,-10,0}, legend pos=south west]
\addplot[gray, dashed, thick] table[x=f, y=dbsk] {datos/bode_modelo.dat};
\addlegendentry{Sallen-Key TP2 (comercial)}
\addplot[only marks, mark=square*, mark size=1.8pt, orange] table[x=f, y=db] {datos/tp2_medido.dat};
\addlegendentry{Sallen-Key TP2 (medido)}
\addplot[black, thin] table[x=f, y=db] {datos/bode_modelo.dat};
\addlegendentry{$H(z)$ identificada}
\addplot[only marks, mark=*, mark size=1.4pt, blue] table[x=f, y=db] {datos/medido_ok.dat};
\addlegendentry{Medido}
\addplot[only marks, mark=x, mark size=2.4pt, red, thick] table[x=f, y=db] {datos/medido_sosp.dat};
\nextgroupplot[ylabel={Fase (°)}, xlabel={$f$ (Hz)}, ymin=-185, ymax=5, ytick={-180,-135,-90,-45,0}]
\addplot[gray, dashed, thick] table[x=f, y=fasesk] {datos/bode_modelo.dat};
\addplot[black, thin] table[x=f, y=fase] {datos/bode_modelo.dat};
\addplot[only marks, mark=*, mark size=1.4pt, blue] table[x=f, y=fase] {datos/medido_ok.dat};
\addplot[only marks, mark=x, mark size=2.4pt, red, thick] table[x=f, y=fase] {datos/medido_sosp.dat};
\end{groupplot}
\end{tikzpicture}
\caption{Respuesta en frecuencia medida, $H(z)$ identificada y filtro Sallen-Key del TP2.}
\label{fig:bode}
\end{figure}

\subsection{Comparación con el Sallen-Key del TP2}
Con los valores comerciales ($R_1 = \SI{8,2}{\kilo\ohm}$, $R_2 = \SI{15}{\kilo\ohm}$, $C_1 = \SI{2,2}{\micro\farad}$, $C_2 = \SI{1}{\micro\farad}$), el Sallen-Key tiene $f_0 = \SI{\resFoSK}{\hertz}$ y $Q = \num{\resQSK}$, de modo que su corte queda en \SI{\resFtresSK}{\hertz} y atenúa \SI{\resDbCincuentaSK}{\dB} a \SI{50}{\hertz}. En el TP2 se midieron \SI{\resTPdosFc}{\hertz} y \SI{\resTPdosDb}{\dB}~\cite{tp2}. El filtro digital, en cambio, tiene el corte en \SI{10}{\hertz}. En el analógico, $f_c$ depende del redondeo a la serie E24 y de la tolerancia de los capacitores (\SI{10}{\percent}), que explica la diferencia de un \SI{12}{\percent} medida en el TP2. En el digital depende solo de coeficientes exactos y del cristal de la placa. En la banda de interés las dos curvas son casi iguales (Fig.~\ref{fig:bode}). Más arriba difieren por la bilineal: el digital atenúa más a medida que se acerca a $f_s/2 = \SI{500}{\hertz}$, donde tiene un cero, y su respuesta es periódica en $f_s$. El Sallen-Key, en cambio, sigue cayendo a \SI{-40}{\dB} por década a cualquier frecuencia, y eso es lo que el digital no puede reproducir (Sección~\ref{sec:aliasing}).

\subsection{Aliasing}
\label{sec:aliasing}
Se aplicaron senoidales de 950, 990 y \SI{1010}{\hertz} (Tabla~\ref{tab:aliasing}). Al muestrear a \SI{1}{\kilo\hertz}, una senoidal de frecuencia $f$ da las mismas muestras que una de $|f - k f_s|$: las tres aparecen en la entrada y en la salida a 50, 10 y \SI{10}{\hertz}. Las pequeñas diferencias (por ejemplo, \SI{50,14}{\hertz}) se deben a que la frecuencia del generador difiere en unos \SI{150}{ppm} del reloj de la placa, como en el barrido. Como el filtro solo ve esas muestras, las trata como señales de 50 y \SI{10}{\hertz}: la de \SI{950}{\hertz} sale atenuada lo mismo que una de \SI{50}{\hertz} (\SI{-28}{\dB}), y las de 990 y \SI{1010}{\hertz} solo unos \SI{3}{\dB}, aunque un pasa bajos de \SI{10}{\hertz} debería eliminarlas. En los tres casos, la ganancia medida coincide dentro de \SI{0,01}{\dB} con $|H|$ en la frecuencia de alias. A esta cadena le falta un filtro analógico antialias antes del ADC, que atenúe todo lo que esté por encima de $f_s/2$. El módulo real sí lo tendría: el propio Sallen-Key, u otro pasa bajos analógico, delante del conversor.

\begin{table}[!htb]
\centering
\caption{Señales por encima de $f_s/2$: frecuencia de alias y ganancia.}
\label{tab:aliasing}
\begin{tabular}{ccccc}
\toprule
$f_\mathrm{gen}$ & $|f - f_s|$ & $f$ medida & $|H|$ medido & $|H(f_\mathrm{alias})|$ \\
(Hz) & (Hz) & (Hz) & (dB) & (dB) \\
\midrule
950 & \num{50} & \num{50,14} & \num{-28,16} & \num{-28,15} \\
990 & \num{10} & \num{10,14} & \num{-3,14} & \num{-3,14} \\
1010 & \num{10} & \num{9,86} & \num{-2,89} & \num{-2,89} \\
\bottomrule

\end{tabular}
\end{table}

\section{Ejercicio 2: pasa banda digital}

\subsection{Especificación, $H(s)$ y frecuencia de muestreo}
El filtro del TP2 que no se implementó es el pasa banda Butterworth de orden~2: $f_0 = \SI{28}{\kilo\hertz}$, $BW = f_{p2} - f_{p1} = \SI{1,7}{\kilo\hertz}$ ($f_{p1} = \SI{\ejFpuno}{\hertz}$, $f_{p2} = \SI{\ejFpdos}{\hertz}$), $A_\mathrm{max} = \SI{0,1}{\dB}$ y ganancia $K = 9$ (\SI{\ejGdb}{\dB}). En el TP2 se lo obtuvo de un prototipo de orden~1 con $\varepsilon = \sqrt{10^{A_\mathrm{max}/10} - 1} = \num{\ejEps}$, como una sola celda
\begin{equation}
H(s) = \frac{K\,(\omega_0/Q)\,s}{s^2 + (\omega_0/Q)\,s + \omega_0^2}, \qquad Q = \varepsilon\frac{f_0}{BW} = \num{\ejQ}.
\label{eq:hspba}
\end{equation}
Se eligió $f_s = \SI{\ejFs}{\kilo\hertz}$. Así, la banda de paso y el tono de rechazo de \SI{56}{\kilo\hertz} de la señal de prueba quedan por debajo de $f_s/2 = \SI{100}{\kilo\hertz}$, con unas 7 muestras por período en $f_0$. El ADC del STM32 llega a 1~Msps y el graficador captura hasta \SI{500}{\kilo\hertz}, pero filtrar en tiempo real a \SI{200}{\kilo\hertz} deja \ejCiclos{} ciclos de \SI{72}{\mega\hertz} por muestra. La rutina del Ejercicio~1 usa 932, lo que limitaría $f_s$ a unos \SI{\ejFsMaxPlaca}{\kilo\hertz}, por lo que un sistema real necesitaría aritmética de punto fijo.

\subsection{Transformación bilineal con predistorsión}
Reemplazando $s = k(1 - z^{-1})/(1 + z^{-1})$, con $k = 2\pi f_0/\tan(\pi f_0/f_s) = \num{\ejKbil}\,\si{\radian\per\second}$, en una sección $(c_2 s^2 + c_1 s + c_0)$ y multiplicando por $(1 + z^{-1})^2$, se obtienen los coeficientes
\begin{equation}
C_0 = c_2 k^2 + c_1 k + c_0, \;\; C_1 = 2(c_0 - c_2 k^2), \;\; C_2 = c_2 k^2 - c_1 k + c_0,
\end{equation}
que \texttt{bilineal\_prewarp()} normaliza por el $A_0$ del denominador. El resultado coincide con el de SciPy, y es
\begin{equation}
H(z) = \frac{b_0\,(1 - z^{-2})}{1 + a_1 z^{-1} + a_2 z^{-2}},
\label{eq:hzpba}
\end{equation}
con $b_0 = \num{\ejBfo}$, $a_1 = \num{\ejAfo}$ y $a_2 = \num{\ejAdfo}$.

\subsection{Cumplimiento de la especificación}
La Tabla~\ref{tab:cumple} muestra que~\eqref{eq:hzpba} no cumple: tiene la ganancia correcta en $f_0$, pero los bordes caen unos \SI{0,13}{\dB}. Cerca de $f_0$ la bilineal comprime la escala de frecuencias en un factor $2\theta/\sin 2\theta$, con $\theta = \pi f_0/f_s$, por lo que la banda digital resulta un \SI{\ejFactor}{\percent} más angosta que la analógica. Subir $f_s$ solo achica el problema: a \SI{500}{\kilo\hertz} los bordes siguen en \SI{-0,104}{\dB}. La corrección adoptada (opción \texttt{--bordes} del script) es predistorsionar los dos bordes: se diseña el pasa banda analógico con bordes en $\Omega_i = k\tan(\pi f_{pi}/f_s)$, $\omega_0 = \sqrt{\Omega_1\Omega_2}$ y $Q = \varepsilon\,\omega_0/(\Omega_2 - \Omega_1)$. Se usa el $\varepsilon$ de \SI{0,09}{\dB}, para dejar \SI{0,01}{\dB} de margen a la cuantización. Así los bordes digitales quedan en $f_{p1}$ y $f_{p2}$ exactos, y la ganancia en $f_0$ sigue siendo \SI{\ejDbCentroB}{\dB}.

\begin{table}[!htb]
\centering
\caption{Ganancia en los bordes respecto de $K$ (especificación: $\geq \SI{-0,1}{\dB}$).}
\label{tab:cumple}
\footnotesize
\setlength{\tabcolsep}{3pt}
\begin{tabular}{lcccccc}
\toprule
Diseño & $f_0$ (Hz) & $Q$ & $b_0$ & $f_{p1}$ (dB) & $f_{p2}$ (dB) & Cumple \\
\midrule
$H(s)$ TP2 & 28000 & \num{\ejQ} & -- & \num{\ejDbSUno} & \num{\ejDbSDos} & Sí \\
Pred.\ $f_0$ & 28000 & \num{\ejQ} & \num{\ejBfo} & \num{\ejDbFoUno} & \num{\ejDbFoDos} & No \\
Pred.\ bordes & \ejFoB & \num{\ejQB} & \num{\ejBb} & \num{\ejDbBUno} & \num{\ejDbBDos} & Sí \\
\bottomrule
\end{tabular}
\end{table}

\subsection{Diagrama de Bode}
En la Fig.~\ref{fig:bodepba}, $H(s)$ y $H(z)$ coinciden cerca de $f_0$, pero se separan a medida que $f$ se acerca a $f_s/2$. $H(z)$ tiene un cero en $z = -1$, así que su módulo se anula en \SI{100}{\kilo\hertz}: a \SI{90}{\kilo\hertz} vale \SI{\ejDbNoventaZ}{\dB}, contra \SI{\ejDbNoventaS}{\dB} de $H(s)$. Por la misma compresión de frecuencias, la banda a \SI{-3}{\dB} pasa de \SIrange{\ejTresSuno}{\ejTresSdos}{\kilo\hertz} en $H(s)$ a \SIrange{\ejTresFouno}{\ejTresFodos}{\kilo\hertz} en $H(z)$, y la corrección la lleva a \SIrange{\ejTresBuno}{\ejTresBdos}{\kilo\hertz}. La fase sigue el mismo patrón.

\begin{figure}[!htb]
\centering
\begin{tikzpicture}
\begin{groupplot}[
    group style={group size=1 by 3, vertical sep=9mm},
    width=\columnwidth, height=3.7cm,
    grid=both, grid style={line width=0.1pt, draw=gray!30}, major grid style={line width=0.2pt, draw=gray!50},
    tick label style={font=\footnotesize}, label style={font=\footnotesize},
    legend style={font=\scriptsize, cells={anchor=west}}, /pgf/number format/use comma]
\nextgroupplot[xmode=log, xmin=1, xmax=100, ymin=-25, ymax=22, ylabel={$|H|$ (dB)}, xlabel={$f$ (kHz)},
    legend pos=south west]
\addplot[gray, dashed, thick] table[x=f, y=hs] {datos/ej2_bode.dat};
\addlegendentry{$H(s)$ TP2}
\addplot[blue, thin] table[x=f, y=hz] {datos/ej2_bode.dat};
\addlegendentry{$H(z)$, pred.\ $f_0$}
\addplot[red, thin] table[x=f, y=hb] {datos/ej2_bode.dat};
\addlegendentry{$H(z)$, pred.\ bordes}
\nextgroupplot[xmin=25, xmax=31, ymin=18.5, ymax=19.2, ylabel={$|H|$ (dB)}, xlabel={$f$ (kHz)},
    yticklabel style={/pgf/number format/fixed, /pgf/number format/precision=1}]
\addplot[gray, dashed, thick] table[x=f, y=hs] {datos/ej2_bode_zoom.dat};
\addplot[blue, thin] table[x=f, y=hz] {datos/ej2_bode_zoom.dat};
\addplot[red, thin] table[x=f, y=hb] {datos/ej2_bode_zoom.dat};
\addplot[black, dotted] coordinates {(25,18.985) (31,18.985)};
\addplot[black, dotted] coordinates {(27.163,18.5) (27.163,19.2)};
\addplot[black, dotted] coordinates {(28.863,18.5) (28.863,19.2)};
\nextgroupplot[xmode=log, xmin=1, xmax=100, ymin=-100, ymax=100, ytick={-90,-45,0,45,90},
    ylabel={Fase (°)}, xlabel={$f$ (kHz)}]
\addplot[gray, dashed, thick] table[x=f, y=fs] {datos/ej2_bode.dat};
\addplot[blue, thin] table[x=f, y=fz] {datos/ej2_bode.dat};
\addplot[red, thin] table[x=f, y=fb] {datos/ej2_bode.dat};
\end{groupplot}
\end{tikzpicture}
\caption{Bode de $H(s)$ y de las dos $H(z)$, con el detalle de la banda de paso: las líneas punteadas marcan $f_{p1}$, $f_{p2}$ y el límite de \SI{-0,1}{\dB}.}
\label{fig:bodepba}
\end{figure}

\subsection{Polos y ceros}
Los ceros están en $z = 1$ y $z = -1$, es decir, en continua y en $f_s/2$. Los polos complejos conjugados están en $|p| = \num{\ejRFo}$ (en el diseño corregido, \num{\ejRB}, a $\pm\num{\ejAngB}^\circ$), dentro del círculo unitario, así que el filtro es estable (Fig.~\ref{fig:pbatiempo}). Para la bilineal con predistorsión en $f_0$, el radio cumple
\begin{equation}
|p|^2 = \frac{1 - x}{1 + x}, \qquad x = \frac{\sin(2\pi f_0/f_s)}{2Q},
\end{equation}
que para $Q = \num{\ejQ}$ da \num{\ejRFormula}, igual al calculado. Cuanto mayor es $Q$, menor es $x$ y más se acercan los polos al círculo. Con $Q \approx 2{,}5$ quedan lejos, lo que hace al filtro poco sensible a la cuantización.

\begin{figure*}[!t]
\centering
\begin{tikzpicture}
\begin{groupplot}[
    group style={group size=3 by 1, horizontal sep=14mm},
    height=4.6cm,
    grid=both, grid style={line width=0.1pt, draw=gray!30}, major grid style={line width=0.2pt, draw=gray!50},
    tick label style={font=\footnotesize}, label style={font=\footnotesize},
    legend style={font=\scriptsize, cells={anchor=west}}, /pgf/number format/use comma]
\nextgroupplot[width=5.2cm, axis equal image, xmin=-1.3, xmax=1.3, ymin=-1.3, ymax=1.3,
    xlabel={Re$(z)$}, ylabel={Im$(z)$}, xtick={-1,0,1}, ytick={-1,0,1}, title={\footnotesize Plano $z$}]
\draw[gray] (axis cs:0,0) circle[radius=1];
\addplot[only marks, mark=x, mark size=3pt, blue, thick] table[x=re, y=im, restrict expr to domain={\thisrow{tipo}}{0:0}] {datos/ej2_pz.dat};
\addplot[only marks, mark=x, mark size=3pt, red, thick] table[x=re, y=im, restrict expr to domain={\thisrow{tipo}}{1:1}] {datos/ej2_pz.dat};
\addplot[only marks, mark=o, mark size=3pt, black, thick] table[x=re, y=im, restrict expr to domain={\thisrow{tipo}}{2:2}] {datos/ej2_pz.dat};
\nextgroupplot[width=6.3cm, xmin=0, xmax=60, xlabel={$n$}, ylabel={Amplitud}, title={\footnotesize Impulso y escalón (corregido)},
    legend pos=north east]
\addplot[red, mark=*, mark size=0.8pt, thin] table[x=n, y=imp] {datos/ej2_tiempo.dat};
\addlegendentry{$h[n]$}
\addplot[black, mark=o, mark size=0.8pt, thin] table[x=n, y=esc] {datos/ej2_tiempo.dat};
\addlegendentry{escalón}
\nextgroupplot[width=6.3cm, xmin=0, xmax=100, ymin=-90, ymax=20, xlabel={$f$ (kHz)}, ylabel={dBV pico},
    title={\footnotesize Señal de prueba}, legend pos=south west]
\addplot[gray, thin] table[x=f, y=ex] {datos/ej2_espectro.dat};
\addlegendentry{entrada}
\addplot[red, thin] table[x=f, y=ey] {datos/ej2_espectro.dat};
\addlegendentry{salida}
\end{groupplot}
\end{tikzpicture}
\caption{Ejercicio 2. Izquierda: polos del diseño con predistorsión en $f_0$ (azul) y corregido (rojo), y ceros en $z = \pm 1$. Centro: respuesta al impulso y al escalón del diseño corregido. Derecha: espectro de la señal de prueba (tonos de 14, 28 y \SI{56}{\kilo\hertz} más ruido) a la entrada y a la salida.}
\label{fig:pbatiempo}
\end{figure*}

\subsection{Respuesta temporal}
La respuesta al impulso arranca en $h[0] = b_0 = \num{\ejHcero}$ y oscila a $f_0$ (unas 7 muestras por período) dentro de una envolvente $|p|^n$, que cae al \SI{1}{\percent} en \ejNuno{} muestras (\SI{\ejTUno}{\micro\second}). La respuesta al escalón salta a $b_0$ y oscila hasta anularse, porque el cero en $z = 1$ bloquea la continua.

\subsection{Señal de prueba}
El script filtra tres tonos de \SI{0,3}{\volt} más ruido. Medidas en el espectro (Fig.~\ref{fig:pbatiempo}), las ganancias son \SI{\ejTonoUnoMed}{\dB} a \SI{\ejTonoUno}{\kilo\hertz}, \SI{\ejTonoDosMed}{\dB} a \SI{\ejTonoDos}{\kilo\hertz} y \SI{\ejTonoTresMed}{\dB} a \SI{\ejTonoTres}{\kilo\hertz}, iguales al centésimo a $|H(e^{j\omega})|$ en esas frecuencias (\SI{\ejTonoUnoTeo}{\dB}, \SI{\ejTonoDosTeo}{\dB} y \SI{\ejTonoTresTeo}{\dB}). Respecto del tono central, los de fuera de banda se atenúan unos \SIrange{11}{13}{\dB}, lo que corresponde a las asíntotas de \SI{\pm 20}{\dB} por década de un pasa banda de orden~2.

\subsection{Cuantización de los coeficientes}
La Tabla~\ref{tab:bits} cuantiza los coeficientes en punto fijo con los bits enteros necesarios (Q2.14 con 16 bits). El diseño con predistorsión en $f_0$ no cumple con ninguna cantidad de bits. El corregido cumple con 16, 14 y 12 bits y deja de cumplir con 11. La menor cantidad que mantiene la especificación es entonces de \textbf{\ejMinBits{} bits}. La estabilidad se mantiene incluso con 4 bits, porque los polos están lejos del círculo unitario: lo que limita es la tolerancia de \SI{0,1}{\dB} de la banda de paso.

\begin{table}[!htb]
\centering
\caption{Cuantización: radio de los polos, peor borde y máximo respecto de $K$ (dB).}
\label{tab:bits}
\footnotesize
\setlength{\tabcolsep}{2.6pt}
\begin{tabular}{ccccccccc}
\toprule
& \multicolumn{4}{c}{Predistorsión en $f_0$} & \multicolumn{4}{c}{Predistorsión en bordes} \\
\cmidrule(lr){2-5}\cmidrule(lr){6-9}
Bits & $|p|$ & Borde & Máx. & Cumple & $|p|$ & Borde & Máx. & Cumple \\
\midrule
$\infty$ & \num{0,857} & \num{-0,131} & \num{0,000} & No & \num{0,830} & \num{-0,090} & \num{0,000} & Sí \\
16 & \num{0,857} & \num{-0,132} & \num{-0,001} & No & \num{0,830} & \num{-0,090} & \num{0,000} & Sí \\
14 & \num{0,857} & \num{-0,134} & \num{-0,003} & No & \num{0,830} & \num{-0,087} & \num{0,003} & Sí \\
12 & \num{0,857} & \num{-0,124} & \num{0,007} & No & \num{0,830} & \num{-0,087} & \num{0,003} & Sí \\
11 & \num{0,857} & \num{-0,131} & \num{0,000} & No & \num{0,829} & \num{-0,117} & \num{-0,024} & No \\
10 & \num{0,857} & \num{-0,131} & \num{0,000} & No & \num{0,829} & \num{-0,120} & \num{-0,024} & No \\
8 & \num{0,857} & \num{-0,103} & \num{0,057} & No & \num{0,829} & \num{-0,103} & \num{0,000} & No \\
6 & \num{0,866} & \num{0,234} & \num{0,470} & No & \num{0,829} & \num{-0,365} & \num{-0,195} & No \\
4 & \num{0,866} & \num{-1,230} & \num{0,915} & No & \num{0,866} & \num{0,353} & \num{2,499} & No \\
\bottomrule

\end{tabular}
\end{table}

\section{Conclusiones}

El filtro de la Blue Pill cumple la especificación de la Etapa~2: es un pasa bajos de orden~2 con ganancia de continua \nm{\resGanDC}, corte en \val{\resFTres}{\hertz} y una respuesta que coincide con el Butterworth de \SI{10}{\hertz} discretizado por bilineal con predistorsión. Frente al Sallen-Key del TP2 tiene la ventaja de no depender de tolerancias de componentes, pero responde igual a cualquier señal que, al muestrearse, caiga en su banda de paso. Por eso no puede reemplazar al Sallen-Key sin un filtro analógico antialias delante del ADC.

En el Ejercicio~2, la bilineal con predistorsión en $f_0$ fija la frecuencia central del pasa banda, pero angosta su banda y no alcanza la tolerancia de \SI{0,1}{\dB} en los bordes. Predistorsionando los dos bordes, con un pequeño margen, el filtro cumple a $f_s = \SI{\ejFs}{\kilo\hertz}$ con coeficientes de \ejMinBits{} bits, y sus polos quedan lejos del círculo unitario.


\begin{thebibliography}{99}
\bibitem{oppenheim}
A. V. Oppenheim y R. W. Schafer, \emph{Discrete-Time Signal Processing}, 3.ª ed. Upper Saddle River, NJ, EE.~UU.: Pearson, 2010, cap.~7.

\bibitem{ieee1057}
\emph{IEEE Standard for Digitizing Waveform Recorders}, IEEE Std 1057-2017, 2018.

\bibitem{tp2}
P. Ernst, V. Rao e I. Gil, \emph{Trabajo Práctico N.º 2: Diseño de filtros activos}, Teoría de Circuitos II, UTN FRC, 2026.

\bibitem{ds}
STMicroelectronics, \emph{STM32F103x8, STM32F103xB: Medium-density performance line Arm-based 32-bit MCU}, hoja de datos DS5319.
\end{thebibliography}

\end{document}
